\begin{definition}[Negation of Continuity at a Point]
\label{def:negation-continuity-at-a-point}
Let \(E \subseteq \mathbb{R}\), let \(f : E \to \mathbb{R}\), and let \(x_0 \in E\).
The statement
\[
f \text{ is not continuous at } x_0
\]
means that there exists \(\varepsilon_0>0\) such that for every \(\delta>0\), there exists a point \(x\in E\) satisfying
\[
|x-x_0|<\delta
\qquad\text{and}\qquad
|f(x)-f(x_0)|\ge \varepsilon_0.
\]
Equivalently,
\[
\neg\Bigl(
\forall \varepsilon>0\;\exists \delta>0\;\forall x\in E:
|x-x_0|<\delta \Rightarrow |f(x)-f(x_0)|<\varepsilon
\Bigr)
\]
is
\[
\exists \varepsilon_0>0\;\forall \delta>0\;\exists x\in E:
\bigl(|x-x_0|<\delta \ \wedge\ |f(x)-f(x_0)|\ge \varepsilon_0\bigr).
\]
\end{definition}

\begin{remark*}[Fully quantified form]
Let \(E \subseteq \mathbb{R}\), let \(f : E \to \mathbb{R}\), and let \(x_0\in E\). Then
\[
f \text{ is not continuous at } x_0
\]
if and only if
\[
\exists \varepsilon_0>0\;\forall \delta>0\;\exists x\in E:
\bigl(|x-x_0|<\delta \ \wedge\ |f(x)-f(x_0)|\ge \varepsilon_0\bigr).
\]
\end{remark*}

\begin{remark*}[Flash]
\FlashQ{What is the negation of continuity at a point?}
\FlashA{A function \(f:E\to\mathbb{R}\) is not continuous at \(x_0\in E\) if and only if there exists \(\varepsilon_0>0\) such that for every \(\delta>0\), there exists \(x\in E\) with
\[
|x-x_0|<\delta
\quad\text{and}\quad
|f(x)-f(x_0)|\ge \varepsilon_0.
\]}
\FlashHint{Negate the quantifiers in order: \(\forall\exists\forall\) becomes \(\exists\forall\exists\), and the conclusion \(<\varepsilon\) becomes \(\ge \varepsilon\).}
\FlashImplies{}
\end{remark*}

\begin{remark*}[Intuition]
This says that some fixed output gap \(\varepsilon_0\) survives no matter how small a neighborhood of \(x_0\) is chosen. Continuity fails precisely when the function values cannot be forced into an arbitrarily small band around \(f(x_0)\).
\end{remark*}